\documentclass[11pt,a4paper]{article}

\usepackage[margin=2cm]{geometry}
\usepackage[T1]{fontenc}
\usepackage{helvet}
\renewcommand{\familydefault}{\sfdefault}
\usepackage[table]{xcolor}
\usepackage{array}
\usepackage{tabularx}
\usepackage{booktabs}
\usepackage{enumitem}
\usepackage{titlesec}
\usepackage{fancyhdr}
\usepackage{listings}
\usepackage[hidelinks]{hyperref}

\definecolor{heading}{HTML}{1F4E79}
\definecolor{rowhead}{HTML}{E6EDF5}
\definecolor{codebg}{HTML}{F5F7FA}

\titleformat{\section}{\Large\bfseries\color{heading}}{\thesection}{0.6em}{}
\titleformat{\subsection}{\large\bfseries}{\thesubsection}{0.6em}{}

\pagestyle{fancy}
\fancyhf{}
\renewcommand{\headrulewidth}{0pt}
\fancyfoot[L]{\small\color{gray}NewsRadar: Mid-Term Progress Report}
\fancyfoot[R]{\small\color{gray}\thepage}

\setlength{\parindent}{0pt}
\setlength{\parskip}{0.5em}
\setlist{itemsep=0.2em, topsep=0.2em}
\renewcommand{\arraystretch}{1.35}
\newcolumntype{L}{>{\raggedright\arraybackslash}X}
\newcommand{\head}[1]{\textbf{#1}}
\newcommand{\code}[1]{\texttt{#1}}

\lstset{
  basicstyle=\ttfamily\small,
  backgroundcolor=\color{codebg},
  frame=single, rulecolor=\color{gray!40},
  breaklines=true, columns=fullflexible,
  xleftmargin=0.5em, xrightmargin=0.5em
}

\begin{document}

{\LARGE\bfseries\color{heading} NewsRadar: Mid-Term Progress Report}\par
\vspace{0.3em}
\textbf{Name:} Deepjyoti Kakoti \quad
\textbf{Roll No.:} \rule{3cm}{0.4pt}\par
\textbf{Branch:} \code{DEEPJYOTI} $\cdot$
\textbf{Repository:} \href{https://github.com/debord555/cs699_project}{github.com/debord555/cs699\_project}\par
\textbf{Module:} AI step (location extraction with a local language model)

% ---------------------------------------------------------------------------
\section{Project summary}

\textbf{NewsRadar -- A Geographical Visualization of News Around the Globe} is a
web platform that places news on a world map. Articles are collected
continuously from many news websites, matched to the places they talk about,
and shown as clickable points (``blips'') on an interactive map. Clicking a
point lists the articles for that place, with links to the original stories.

\begin{tabularx}{\textwidth}{>{\bfseries}p{4.2cm} L}
\toprule
\rowcolor{rowhead} \head{Feature} & \head{What the user gets} \\
\midrule
Map of news & Every news event appears as a point on a world map, at the place it is about. \\
Continuous collection & News websites are scraped regularly, so the map stays up to date. \\
No duplicate events & Several articles about the same real event are grouped into one point. \\
Categories & News is tagged as Business, Sports, Military, Disasters, Crime, Environment, Technology, Opportunities and more. \\
Time filter & Users can choose a date range to see recent or older events. \\
Zoom & Users can zoom into any region and explore the news there. \\
Text search & A search such as ``earthquake'' shows matching news from around the world on the map. \\
Original sources & Every article links back to the website that published it. \\
\bottomrule
\end{tabularx}

% ---------------------------------------------------------------------------
\section{My part: the AI step}

To put an article on the map, the system has to know \emph{where} the story
happens. Articles arrive as plain text, so a language model reads each one and
lists the places it mentions, with their latitude and longitude. My work in
this phase was to:

\begin{enumerate}
  \item set up a language model that runs on our own computer, with one script; and
  \item write the worker that sends each article to the model and stores the places it finds in the database.
\end{enumerate}

The worker sits between the scraper, which saves articles, and the map, which
reads places from the database.

% ---------------------------------------------------------------------------
\section{Technology choices}

\begin{tabularx}{\textwidth}{>{\bfseries}p{3.2cm} p{4.3cm} L}
\toprule
\rowcolor{rowhead} \head{Choice} & \head{What it is} & \head{Why it was chosen} \\
\midrule
Lemonade Server & A free program that runs language models on a local PC & No cost per request, no internet needed, and article text never leaves our machine. It offers the same interface (API) as OpenAI, so the code would work with a hosted model too. \\
Qwen3-8B (Q4\_K\_M) & An open language model with 8 billion parameters, compressed to about 5\,GB & Small enough to run on a normal GPU, yet good at reading text and returning structured answers such as JSON. \\
Python & A programming language & Same language as the rest of the project; ready-made libraries for databases (\code{psycopg}) and web requests (\code{requests}). \\
PowerShell & The Windows command shell & Lets the whole model setup run as one script on our Windows machines. \\
\bottomrule
\end{tabularx}

% ---------------------------------------------------------------------------
\section{Setting up the model}

The script \code{scripts/setup\_llm.ps1} prepares the model in one command:

\begin{lstlisting}
powershell -ExecutionPolicy Bypass -File scripts\setup_llm.ps1
\end{lstlisting}

It runs these steps in order and stops with a clear message if one fails:

\begin{enumerate}
  \item \textbf{Find Lemonade.} Looks for the \code{lemonade} program, first on the system path and then in its default install folder.
  \item \textbf{Start the server.} If the server is not already running, starts it in the background and waits up to 60 seconds for its health check to pass.
  \item \textbf{Download the model.} Downloads Qwen3-8B (about 5\,GB) only if it is not already on the machine.
  \item \textbf{Load the model} into memory.
  \item \textbf{Smoke test.} Asks the model to reply with ``OK'' to confirm the whole chain works.
  \item \textbf{Print the settings} (\code{LLM\_BASE\_URL}, \code{LLM\_MODEL}) to copy into the private \code{.env} file.
\end{enumerate}

The script is safe to run again: steps that are already done are skipped. The
model name and port can be changed with the \code{-Model} and \code{-Port}
options. I also added the two LLM settings to \code{.env.example}, so every
teammate knows which values to set.

% ---------------------------------------------------------------------------
\section{The location extractor}\label{sec:locate}

\subsection{How it works}

\code{llm/locate.py} runs as a loop that never stops:

\begin{enumerate}
  \item Look for one article that is waiting to be processed: status \code{ready} and full text present. If there is none, wait 2 seconds and look again.
  \item Send the article's text to the model and ask for the places it mentions.
  \item For each place, check whether it is already in the \code{locations} table (ignoring upper/lower case). If not, add it with its coordinates.
  \item Link the article to each place in the \code{article\_location\_map} table, keeping the order in which the places were mentioned.
  \item Mark the article \code{llm\_done}. If the model call failed, mark it \code{llm\_failed} and save the error message instead.
  \item Go back to step 1.
\end{enumerate}

\subsection{What the model is asked}

The model is told to reply with \emph{only} a JSON list, so the program can read
the answer directly:

\begin{lstlisting}
[{"name": "Guwahati", "latitude": 26.14, "longitude": 91.73},
 {"name": "Assam",    "latitude": 26.20, "longitude": 92.94}]
\end{lstlisting}

The temperature is set to 0, so the same article always gives the same answer.
The prompt ends with \code{/no\_think}, which turns off Qwen3's long
``thinking'' output and makes replies much faster.

\subsection{Key design decisions}

\begin{tabularx}{\textwidth}{>{\bfseries}p{5.2cm} L}
\toprule
\rowcolor{rowhead} \head{Decision} & \head{Reason} \\
\midrule
Use the article \code{status} column to find work & An article with no places would never get a row in \code{article\_location\_map}, so checking that table would pick the same article again forever. The status (\code{ready} $\rightarrow$ \code{llm\_done} / \code{llm\_failed}) marks every article as finished exactly once. \\
One database transaction per article & An article's places and its new status are saved together. If the program crashes halfway, nothing is saved and the article is simply tried again. \\
Lock the article while working on it (\code{FOR UPDATE SKIP LOCKED}) & Several copies of the worker can run side by side without two of them taking the same article. \\
Clean the model's answer before saving & Small models sometimes add extra text or repeat a place. The worker keeps only the JSON list, drops repeated names, and discards coordinates outside the valid range (latitude $-90$ to $90$, longitude $-180$ to $180$). \\
Reuse existing places & ``Guwahati'' and ``guwahati'' point to the same row, so the map shows one point per place rather than many. \\
Cut very long articles to 12{,}000 characters & Keeps each request within what the model can read at once. \\
Failed articles are marked, not retried forever & One bad article cannot block the queue, and the saved error explains what went wrong. \\
\bottomrule
\end{tabularx}

\subsection{Settings}

\begin{tabularx}{\textwidth}{>{\ttfamily}p{4.5cm} p{3cm} L}
\toprule
\rowcolor{rowhead} \head{\sffamily Setting} & \head{Value} & \head{Meaning} \\
\midrule
LLM\_BASE\_URL & \code{127.0.0.1:13305} & Address of the local model server \\
LLM\_MODEL & Qwen3-8B & Model to use \\
LLM\_TIMEOUT & 300 seconds & Longest wait for one reply from the model \\
LLM\_MAX\_BODY\_CHARS & 12{,}000 & Longest article text sent to the model \\
LOCATE\_IDLE\_SECONDS & 2 seconds & Pause when there is no work \\
\bottomrule
\end{tabularx}

% ---------------------------------------------------------------------------
\section{Description of the files}

\begin{tabularx}{\textwidth}{>{\ttfamily}p{4.6cm} L}
\toprule
\rowcolor{rowhead} \head{\sffamily File} & \head{Purpose} \\
\midrule
scripts/setup\_llm.ps1 & Installs, starts, loads and tests the local language model with one command. \\
.env.example & Template settings; I added \code{LLM\_BASE\_URL} and \code{LLM\_MODEL}. \\
llm/locate.py & The location extractor described in Section~\ref{sec:locate}. \\
llm/db.py & Opens a connection to the database using the private \code{.env} settings. \\
llm/settings.py & Timeouts and limits for the extractor. \\
llm/requirements.txt & The three Python libraries the extractor needs, with exact versions. \\
\bottomrule
\end{tabularx}

To run the extractor:

\begin{lstlisting}
cd llm
pip install -r requirements.txt
python locate.py
\end{lstlisting}

% ---------------------------------------------------------------------------
\section{Outcome of this phase}

\begin{itemize}
  \item A local language model that any teammate can set up and test with one command.
  \item A working location extractor that turns article text into places with coordinates and links them to the article in the database.
  \item Safe for several workers at once, and never leaves half-saved results.
  \item All work is on GitHub (branch \code{DEEPJYOTI}).
\end{itemize}

% ---------------------------------------------------------------------------
\section{Next steps}

\begin{itemize}
  \item Run the extractor on real articles from the scraper and measure how often the places and coordinates are correct.
  \item Extend the same step to tag each article with a news category.
\end{itemize}

\end{document}
